% Style based on the supplied Waseem_Generic_AI_ML_GenAI_AgenticAI.tex.
% Python replaces template tokens. Edit layout here, content in JSON/portfolio JS.
\documentclass[9pt,letterpaper]{extarticle}
\usepackage[left=0.68in,right=0.68in,top=0.45in,bottom=0.45in]{geometry}
\usepackage[T1]{fontenc}
\input{glyphtounicode}
\pdfgentounicode=1
\usepackage[utf8]{inputenc}
\usepackage{lmodern}
\usepackage{microtype}
\usepackage{pifont}
\usepackage{xcolor}
\usepackage{enumitem}
\usepackage{titlesec}
\usepackage[hidelinks]{hyperref}
\usepackage{tabularx,array,parskip,ragged2e,needspace}
\pagestyle{empty}
\setlength{\parindent}{0pt}
\setlength{\parskip}{1pt}
\definecolor{accent}{HTML}{155B8A}
\titleformat{\section}{\large\bfseries\color{accent}}{}{0pt}{}[\vspace{0.5pt}\color{accent}\titlerule]
\titlespacing*{\section}{0pt}{3pt}{2pt}
% Compact subsections beneath the main accent-rule sections.
\titleformat{\subsection}{\normalsize\bfseries\color{accent}}{}{0pt}{}
\titlespacing*{\subsection}{0pt}{3pt}{1pt}
\setlist[itemize]{leftmargin=1.15em,itemsep=0.8pt,topsep=1pt,parsep=0pt,partopsep=0pt}
\newcommand{\role}[4]{%
  \Needspace{4\baselineskip}\textbf{#1}\hfill\textbf{#2}\\[-1pt]
  \textit{#3}\hfill\textit{#4}\par}
\newcommand{\entry}[2]{\Needspace{4\baselineskip}\textbf{#1}\hfill\textit{#2}\par}
\newcommand{\skillrow}[2]{\textbf{#1:} & #2\\}
\hypersetup{pdftitle={Waseem Raza, PhD - Faculty / Generic},pdfauthor={Waseem Raza, PhD}}
\begin{document}
\hyphenpenalty=10000\exhyphenpenalty=10000\pretolerance=10000\emergencystretch=2em\widowpenalty=10000\clubpenalty=10000
\begin{center}
{\LARGE\bfseries Waseem Raza, PhD}{\normalsize\mdseries \enspace (Electrical and Computer Engineering)}\\[2pt]
{\large\bfseries Researcher \& Engineering Educator \textbar{} AI/ML for Wireless Networks}\\[1pt]
{\small Wireless AI \textbar{} 5G/6G Networks \textbar{} Robust Learning \textbar{} Research \& Teaching}\\[3pt]
{\footnotesize \ding{41}\enspace \href{mailto:waseem.placements@gmail.com}{waseem.placements@gmail.com} \textbar{} \href{mailto:waseemraza844@gmail.com}{waseemraza844@gmail.com}\\[1pt]New Jersey, United States \quad \ding{37}\enspace \href{tel:+17328039703}{+1 732-803-9703}\\[1pt]\href{https://www.linkedin.com/in/waseem-raza-phd-9a737858/}{LinkedIn} \textbar{} \href{https://github.com/WaseemRaza844}{GitHub} \textbar{} \href{https://scholar.google.com/citations?user=d58Drr0AAAAJ\&hl=en}{Google Scholar} \textbar{} \href{https://waseemraza844.github.io/portfolio/learning.html}{Coursera Courses} \textbar{} \href{https://waseemraza844.github.io/portfolio/}{Portfolio}}
\end{center}
\vspace{-5pt}
\Needspace{3\baselineskip}
\section{Academic Profile}
Electrical and Computer Engineering researcher and educator with a PhD from the University of Oklahoma specializing in AI/ML for wireless networks. Research addresses reliable learning, propagation modeling, network optimization, self-healing, digital twins, generative augmentation, and reinforcement learning for 5G/6G systems. University teaching and research are complemented by industrial R\&D at AT\&T Labs, Futurewei, and AT\&T, bridging rigorous experiments with real-world network challenges.
\Needspace{2\baselineskip}
\subsection{Research Interests}
Wireless AI and 5G/6G; Robust and Data-Efficient ML; Digital Twins and Propagation Modeling; Network Optimization and Self-Healing; Transfer and Generative Learning; Autonomous Networks and Wearable AI\par
\Needspace{2\baselineskip}
\subsection{Doctoral Research Focus}
Doctoral research in AI/ML for wireless networks, emphasizing resilient propagation models, data-efficient generative learning, adaptive network optimization, self-healing, and reproducible evaluation.\par
\Needspace{3\baselineskip}
\section{Education}
\entry{PhD, Electrical and Computer Engineering}{09/2019 - 07/2025}
University of Oklahoma, Norman, Oklahoma, USA\par
\entry{MSc, Telecommunication Engineering}{08/2014 - 08/2016}
University of Engineering and Technology, Taxila, Pakistan\par
\entry{BSc, Telecommunication Engineering}{11/2010 - 06/2014}
University of Engineering and Technology, Taxila, Pakistan\par
\Needspace{3\baselineskip}
\section{Teaching, Research \& Professional Experience}
\Needspace{6\baselineskip}
\Needspace{2\baselineskip}
\subsection{Teaching \& Instruction}
\Needspace{6\baselineskip}
\role{Lecturer, Department of Technology}{12/2017 - 09/2019}{University of Lahore, Pakistan}{Lahore, Pakistan}
\begin{itemize}\interlinepenalty=10000
\item Taught programming, communication systems, and wireless communication with practical examples.
\item Advised senior-design students on projects, academic plans, and early careers.
\item Aligned course materials with learning outcomes and accreditation standards.
\item Supported accreditation, exams, technical workshops, and student activities.
\end{itemize}
\vspace{1pt}
{\small\textbf{Teaching areas:} Wireless Communication Systems; LTE/5G Networks; Signals and Digital Signal Processing; Machine Learning and AI; Programming and Data Science; Experimental Design and Labs}\par
\Needspace{6\baselineskip}
\Needspace{2\baselineskip}
\subsection{University Research Appointments}
\Needspace{6\baselineskip}
\role{Graduate Research Assistant}{09/2019 - 07/2025}{AI4Networks, University of Oklahoma}{Tulsa, OK}
\begin{itemize}\interlinepenalty=10000
\item Built 3GPP-compliant Python simulation tools and an AI-enabled 5G research testbed.
\item Built GAN, transfer-learning, autoencoder, and reinforcement-learning models for wireless data.
\item Led 100+ experiments testing wireless ML under domain shift and limited labeled data.
\item Developed self-healing and outage-management methods, improving resilience by 30\%.
\item Published wireless AI research and mentored junior researchers in experimental design.
\end{itemize}
\vspace{1pt}
\Needspace{6\baselineskip}
\role{Research Associate, ComSens Lab}{06/2016 - 12/2017}{COMSATS University Islamabad, Pakistan}{Islamabad, Pakistan}
\begin{itemize}\interlinepenalty=10000
\item Modeled wireless communication systems through simulation, analysis, and technical reports.
\item Mentored Bachelor's and Master's theses in experiment design and research presentation.
\item Reviewed student reports and manuscripts for rigor and publication readiness.
\item Supported proposals, literature reviews, lab seminars, and research coordination.
\end{itemize}
\vspace{1pt}
\Needspace{6\baselineskip}
\role{Full-Time Research Scholar, Communication Systems Lab}{08/2014 - 02/2016}{University of Engineering and Technology, Taxila, Pakistan}{Taxila, Pakistan}
\begin{itemize}\interlinepenalty=10000
\item Researched wireless communication systems while completing an MSc in Telecommunications.
\item Supported communications courses, labs, assignments, grading, and examinations.
\item Led wireless communication labs and prepared technical manuals and experiments.
\item Developed course and student-project materials for practical wireless engineering.
\end{itemize}
\vspace{1pt}
\Needspace{6\baselineskip}
\Needspace{2\baselineskip}
\subsection{Industry Research \& Engineering}
\Needspace{6\baselineskip}
\role{Machine Learning Engineer}{02/2025 - Present}{AT\&T, contractor via Innovatix Inc.}{Bedminster, NJ}
\begin{itemize}\interlinepenalty=10000
\item Applied domain-adaptive ML to telecom service assurance and model validation.
\item Designed cross-domain evaluation and error-analysis methods for operational ML.
\item Worked with engineering and public-safety teams on reliable network assurance.
\end{itemize}
\vspace{1pt}
\Needspace{6\baselineskip}
\role{Advanced-Degrees Intern}{06/2024 - 08/2024}{AT\&T Labs}{Austin, TX}
\begin{itemize}\interlinepenalty=10000
\item Analyzed 500K+ LTE/5G-NR logs with ML, reducing diagnostic time by 20\%.
\item Improved 5G NR-NTN cell selection by 15\% through data-driven validation.
\item Presented operational data insights to next-generation wireless engineering teams.
\end{itemize}
\vspace{1pt}
\Needspace{6\baselineskip}
\role{Advanced Wireless Technology Development Intern}{09/2022 - 01/2023}{Futurewei Technologies}{Schaumburg, IL}
\begin{itemize}\interlinepenalty=10000
\item Improved wireless CSI compression by 25\% with transfer learning and autoencoders.
\item Ran 50+ experiments on wireless ML robustness under changing data conditions.
\item Translated research results into recommendations for future wireless networks.
\end{itemize}
\vspace{1pt}
\Needspace{6\baselineskip}
\role{Advanced-Degrees Intern}{06/2022 - 08/2022}{AT\&T Labs}{Bedminster, NJ}
\begin{itemize}\interlinepenalty=10000
\item Improved 5G standalone user-timing precision by 20\% with ML and statistics.
\item Analyzed user traces with SQL and presented results to engineering leadership.
\item Turned operational traces into indicators for network diagnosis and optimization.
\end{itemize}
\vspace{1pt}
\Needspace{3\baselineskip}
\section{Research Contributions}
\Needspace{2\baselineskip}
\subsection{Publications}
\Needspace{4\baselineskip}\noindent\begin{tabularx}{\linewidth}{@{}>{\raggedright\arraybackslash}p{0.68\linewidth}@{\hspace{8pt}}>{\raggedleft\arraybackslash}X@{}}
\textbf{Real-World Medication Adherence Monitoring Using Hybrid Neural Networks and Smartwatches} & {\scriptsize\textit{HealthINF \textbar{} Accepted}\newline 2026}\\
\end{tabularx}\vspace{1pt}
\Needspace{4\baselineskip}\noindent\begin{tabularx}{\linewidth}{@{}>{\raggedright\arraybackslash}p{0.68\linewidth}@{\hspace{8pt}}>{\raggedleft\arraybackslash}X@{}}
\textbf{\href{https://doi.org/10.1016/j.comcom.2025.108129}{AI-Powered Resilience: A Dual-Approach for Outage Management in Dense Cellular Networks}} & {\scriptsize\textit{Computer Communications}\newline 2025 \textbar{} \href{https://doi.org/10.1016/j.comcom.2025.108129}{DOI}}\\
\end{tabularx}\vspace{1pt}
\Needspace{4\baselineskip}\noindent\begin{tabularx}{\linewidth}{@{}>{\raggedright\arraybackslash}p{0.68\linewidth}@{\hspace{8pt}}>{\raggedleft\arraybackslash}X@{}}
\textbf{\href{https://doi.org/10.3390/s25247651}{On Multi-Parameter Optimization and Proactive Reliability in 5G and Beyond Cellular Networks}} & {\scriptsize\textit{Sensors 2025}\newline 2025 \textbar{} \href{https://doi.org/10.3390/s25247651}{DOI}}\\
\end{tabularx}\vspace{1pt}
\Needspace{4\baselineskip}\noindent\begin{tabularx}{\linewidth}{@{}>{\raggedright\arraybackslash}p{0.68\linewidth}@{\hspace{8pt}}>{\raggedleft\arraybackslash}X@{}}
\textbf{\href{https://doi.org/10.1109/PIMRC59610.2024.10817403}{An AI-Driven Framework for Enhancing Resilience in Propagation Models to Enable Digital Twin}} & {\scriptsize\textit{IEEE PIMRC}\newline 2024 \textbar{} \href{https://doi.org/10.1109/PIMRC59610.2024.10817403}{DOI}}\\
\end{tabularx}\vspace{1pt}
\Needspace{4\baselineskip}\noindent\begin{tabularx}{\linewidth}{@{}>{\raggedright\arraybackslash}p{0.68\linewidth}@{\hspace{8pt}}>{\raggedleft\arraybackslash}X@{}}
\textbf{\href{https://www.ai4networks.com/files/conferences/c-24-7.pdf}{Refining Wireless Propagation Models using Domain-Informed GANs amid Data Scarcity}} & {\scriptsize\textit{IEEE VTC-Fall 2024}\newline 2024 \textbar{} \href{https://www.ai4networks.com/files/conferences/c-24-7.pdf}{Link}}\\
\end{tabularx}\vspace{1pt}
\Needspace{4\baselineskip}\noindent\begin{tabularx}{\linewidth}{@{}>{\raggedright\arraybackslash}p{0.68\linewidth}@{\hspace{8pt}}>{\raggedleft\arraybackslash}X@{}}
\textbf{\href{https://doi.org/10.1109/MNET.006.2100377}{An AI-Enabled Framework to Defend Ingenious MDT-Based Attacks on the Emerging Zero Touch Cellular Networks}} & {\scriptsize\textit{IEEE Network, 38(1), 228-237}\newline 2024 \textbar{} \href{https://doi.org/10.1109/MNET.006.2100377}{DOI}}\\
\end{tabularx}\vspace{1pt}
\Needspace{4\baselineskip}\noindent\begin{tabularx}{\linewidth}{@{}>{\raggedright\arraybackslash}p{0.68\linewidth}@{\hspace{8pt}}>{\raggedleft\arraybackslash}X@{}}
\textbf{\href{https://arxiv.org/abs/2402.15795}{Positioning Error Impact Compensation through Data-Driven Optimization in User-Centric Networks}} & {\scriptsize\textit{IEEE GLOBECOM 2023}\newline 2023 \textbar{} \href{https://arxiv.org/abs/2402.15795}{Link}}\\
\end{tabularx}\vspace{1pt}
\Needspace{4\baselineskip}\noindent\begin{tabularx}{\linewidth}{@{}>{\raggedright\arraybackslash}p{0.68\linewidth}@{\hspace{8pt}}>{\raggedleft\arraybackslash}X@{}}
\textbf{\href{https://arxiv.org/abs/2203.15899}{Toward a Hybrid RF/Optical Lunar Communication System (LunarComm)}} & {\scriptsize\textit{IEEE Network}\newline 2022 \textbar{} \href{https://arxiv.org/abs/2203.15899}{Link}}\\
\end{tabularx}\vspace{1pt}
\Needspace{2\baselineskip}
\subsection{Research Projects \& Platforms}
\entry{SyntheticNET: 3GPP-Compliant Network Simulation}{Research Project \textbar{} 2020-2024}
Developed a Python-based, 3GPP-aligned simulation environment for realistic mobility, propagation, and cellular-network experimentation when measured data is sparse or costly. Supported controlled data generation and repeatable studies in propagation modeling, network optimization, digital twins, and resilient wireless AI.\par
{\small\textit{Tools \& Methods:} Python, 3GPP, Wireless Simulation, Mobility Modeling, Propagation Modeling, Synthetic Data, Experiment Design, Digital Twins, Network Analytics\par}
\vspace{1pt}
\entry{TurboRAN Experimental Platform}{Research Infrastructure \textbar{} 2020-2024}
Built and used a multi-band, multi-tier experimental platform for evaluating AI-enabled radio-access-network methods under controlled network conditions. Supported repeatable measurement collection, testbed experimentation, RAN optimization studies, and validation of emerging wireless-network algorithms.\par
{\small\textit{Tools \& Methods:} 5G Testbeds, RAN, Network Measurement, Experiment Design, Wireless Prototyping, Multi-Band Systems, Performance Evaluation, RAN Optimization\par}
\vspace{1pt}
\entry{Propagation Models \& Digital Twins}{Research Project \textbar{} 2024}
Developed robust data-driven radio-propagation models for wireless digital twins using domain-informed generative learning and scarce measured data. Combined 3GPP-informed pretraining, GAN-based augmentation, domain adaptation, and statistical validation to improve resilience under distribution shift.\par
{\small\textit{Tools \& Methods:} Propagation Modeling, Digital Twins, GANs, Domain Adaptation, 3GPP Models, Synthetic Data, Statistical Validation, Wireless AI, Distribution Shift\par}
\vspace{1pt}
\entry{Wearable Intelligence for Health Monitoring}{Research Project \textbar{} 2024-2026}
Developed deep-learning workflows for recognizing medication-taking and dietary activities from smartwatch accelerometer and gyroscope signals in free-living settings. Evaluated windowing, feature extraction, CNN/LSTM-style sequence models, cross-user generalization, and recall-oriented performance for health-monitoring applications.\par
{\small\textit{Tools \& Methods:} Wearable AI, Smartwatch IMU, Time Series, Deep Learning, CNN/LSTM, Feature Engineering, Cross-User Evaluation, Activity Recognition, Healthcare AI\par}
\vspace{1pt}
\Needspace{3\baselineskip}
\section{Research Methods \& Technical Expertise}
\Needspace{2\baselineskip}\noindent{\small\textbf{Machine Learning:} Deep learning, transfer learning, domain adaptation, adversarial learning, generative modeling, reinforcement learning, classification, time-series analysis, robustness evaluation, and model validation}\par
\Needspace{2\baselineskip}\noindent{\small\textbf{Wireless \& Networked Systems:} LTE/4G, 5G-NR, NTN, mobility, RAN, propagation modeling, outage management, self-healing, network measurements, and service assurance}\par
\Needspace{2\baselineskip}\noindent{\small\textbf{Research Methods:} Experimental design, 3GPP-aligned simulation, statistical analysis, cross-validation, reproducible workflows, distribution-shift evaluation, technical writing, and publication development}\par
\Needspace{2\baselineskip}\noindent{\small\textbf{Engineering Tools:} Python, PyTorch, TensorFlow, scikit-learn, XGBoost, Pandas, NumPy, SQL, cloud analytics, data pipelines, visualization, Git, and Linux}\par
\Needspace{2\baselineskip}\noindent{\small\textbf{Teaching \& Mentoring:} Wireless communications, communication systems, programming, machine learning, laboratory instruction, experimental methods, student projects, and research development}\par
\Needspace{3\baselineskip}
\section{Academic Service \& Recognition}
\entry{William H. Barkow Scholarship}{2024}
{\small University of Oklahoma, Electrical and Computer Engineering\par}
\entry{Clyde L. Farrar Endowment Scholarship}{2024}
{\small University of Oklahoma\par}
\entry{P.H. Robinson Fellowship}{2023}
{\small University of Oklahoma\par}
\entry{FAS-Charles Forgan Scholarship}{2023}
{\small University of Oklahoma\par}
\entry{Nettie Vincent Boggs Graduate Endowment}{2022}
{\small University of Oklahoma\par}
\Needspace{3\baselineskip}
\section{Self-Learning}
\Needspace{4\baselineskip}\noindent\begin{tabularx}{\linewidth}{@{}>{\raggedright\arraybackslash}p{0.59\linewidth}@{\hspace{7pt}}>{\raggedleft\arraybackslash}X@{}}
\textbf{1. \href{https://www.coursera.org/account/accomplishments/verify/VU3Y5TBLVLZT}{Machine Learning}} & {\scriptsize \href{https://www.coursera.org/account/accomplishments/verify/VU3Y5TBLVLZT}{Stanford/Coursera} \textbar{} Completed \textbar{} Mar 2020}\\
\end{tabularx}\vspace{-2pt}
\noindent{\footnotesize Supervised and unsupervised learning, regularization, model evaluation, and practical predictive modeling.}\par\vspace{2pt}
\Needspace{4\baselineskip}\noindent\begin{tabularx}{\linewidth}{@{}>{\raggedright\arraybackslash}p{0.59\linewidth}@{\hspace{7pt}}>{\raggedleft\arraybackslash}X@{}}
\textbf{2. \href{https://www.coursera.org/account/accomplishments/verify/Y6NP4GWLS4DZ}{Introduction to Deep Learning \& Neural Networks with Keras}} & {\scriptsize \href{https://www.coursera.org/account/accomplishments/verify/Y6NP4GWLS4DZ}{IBM/Coursera} \textbar{} Completed \textbar{} Aug 2026}\\
\end{tabularx}\vspace{-2pt}
\noindent{\footnotesize Neural networks, training workflows, optimization, loss functions, and model evaluation with Keras.}\par\vspace{2pt}
\Needspace{4\baselineskip}\noindent\begin{tabularx}{\linewidth}{@{}>{\raggedright\arraybackslash}p{0.59\linewidth}@{\hspace{7pt}}>{\raggedleft\arraybackslash}X@{}}
\textbf{3. \href{https://waseemraza844.github.io/portfolio/certificates/wireless/c1-ai-found-indus-overview.pdf}{AI Foundations \& Industry Overview for Telecommunication}} & {\scriptsize \href{https://www.coursera.org/account/accomplishments/verify/3900EVOH4I1Y}{Coursera} \textbar{} Completed \textbar{} Jun 2026}\\
\end{tabularx}\vspace{-2pt}
\noindent{\footnotesize AI applications in telecom networks, operations analytics, service assurance, and automation.}\par\vspace{2pt}
\Needspace{4\baselineskip}\noindent\begin{tabularx}{\linewidth}{@{}>{\raggedright\arraybackslash}p{0.59\linewidth}@{\hspace{7pt}}>{\raggedleft\arraybackslash}X@{}}
\textbf{4. \href{https://waseemraza844.github.io/portfolio/certificates/wireless/c2-net-sec-optim-tel.pdf}{Network and Security Optimization in Telecommunication}} & {\scriptsize \href{https://www.coursera.org/account/accomplishments/verify/4TQ51W1OQWRE}{Coursera} \textbar{} Completed \textbar{} Jun 2026}\\
\end{tabularx}\vspace{-2pt}
\noindent{\footnotesize Data-driven network optimization, telecom security, reliability analysis, and operational efficiency.}\par\vspace{2pt}
\Needspace{4\baselineskip}\noindent\begin{tabularx}{\linewidth}{@{}>{\raggedright\arraybackslash}p{0.59\linewidth}@{\hspace{7pt}}>{\raggedleft\arraybackslash}X@{}}
\textbf{5. \href{https://waseemraza844.github.io/portfolio/certificates/wireless/4g-network-fundamentals.pdf}{4G Network Fundamentals}} & {\scriptsize \href{https://www.coursera.org/account/accomplishments/verify/DLW1YN474VZJ}{Mines-Télécom/Coursera} \textbar{} Completed \textbar{} Aug 2026}\\
\end{tabularx}\vspace{-2pt}
\noindent{\footnotesize LTE radio and core networks, cellular architecture, mobility management, and 4G fundamentals.}\par\vspace{2pt}
\Needspace{4\baselineskip}\noindent\begin{tabularx}{\linewidth}{@{}>{\raggedright\arraybackslash}p{0.59\linewidth}@{\hspace{7pt}}>{\raggedleft\arraybackslash}X@{}}
\textbf{6. \href{https://www.coursera.org/learn/5g-network-fundamentals}{5G Network Fundamentals}} & {\scriptsize \href{https://www.coursera.org/learn/5g-network-fundamentals}{Mines-Télécom/Coursera} \textbar{} In progress (\textasciitilde{}60\%)}\\
\end{tabularx}\vspace{-2pt}
\noindent{\footnotesize 5G NR radio, network architecture, service capabilities, and next-generation cellular fundamentals.}\par\vspace{2pt}
\end{document}
